% ===================================================================
%  جدول نمبر 10 — سمندر۔ بندرگاہ۔
%
%  Traduction, faite en 2026, en ourdou d'aujourd'hui, sur l'ido de
%  texto/io. Regles de la colonne : voir 10-jadval-01.tex. Une seule
%  scene, une seule numerotation. Ecrit avec outils/ordo.py sous les
%  yeux.
%
%  LA LIGNE DE PARTAGE DE CE TABLEAU N'EST PAS LE RIVAGE, C'EST LA
%  VAPEUR. Tout ce que la marine a voile connaissait porte en ourdou
%  un nom persan ou arabe, et pas un seul emprunt : آبدوز le
%  sous-marin — « le plongeur d'eau » —, عرشہ le pont, مستول le mat,
%  لنگر l'ancre, بادبان la voile, چپّو la rame, گودی le bassin,
%  بندرگاہ le port, آبنائے le detroit, برزخ l'isthme, جزیرہ نما la
%  presqu'ile, خلیج la baie, مد و جزر la maree, افق l'horizon,
%  محصول خانہ la douane, فصیل le rempart. Tout ce que la vapeur a
%  invente est anglais : کروزر le croiseur (6), فریگیٹ la fregate
%  (11), تارپیڈو la torpille (8), ٹگ le remorqueur (16), پروپیلر
%  l'helice (28), بوائلر la chaudiere (25). Vingt-trois mots d'un
%  cote, treize de l'autre, et la coupure tombe a la meme date pour
%  tous : le milieu du XIXe siecle.
%
%  ET LES GRADES SONT ANGLAIS AUSSI, POUR UNE AUTRE RAISON.
%  لیفٹیننٹ, کمانڈر, ایڈمرل, کوارٹر ماسٹر : la marine du Pakistan
%  est nee de la Royal Indian Navy et en a garde les titres, comme
%  l'armee de terre a garde کرنل et میجر. Le persan avait pourtant
%  ses mots — امیر البحر pour l'amiral, qui est meme l'ancetre du
%  mot « admiral » lui-meme, revenu par l'arabe et l'italien. La
%  colonne ne le remet pas : personne au Pakistan n'appelle un
%  amiral امیر البحر, et une traduction de 2026 ne ressuscite pas un
%  titre pour faire joli. C'est le meme principe qu'au chiffre tamoul
%  du tableau 1 : la langue qu'on ecrit est celle qu'on lit.
%
%  آبدوز MERITE QU'ON S'Y ARRETE. Le mot est un compose persan
%  transparent — آب l'eau, دوز qui plonge — forge au XXe siecle pour
%  une chose que le persan n'avait jamais vue, et il a pris partout
%  sans qu'on ait besoin de « submarine ». C'est l'exception qui
%  eclaire la regle du glacier au tableau 9 : quand une langue forge
%  a temps, l'emprunt ne s'installe pas. گلیشیئر est arrive avant que
%  l'ourdou ait forge ; آبدوز est arrive avant l'anglais.
%
%  LA PREDICTION DU TABLEAU 9 SE VERIFIE. On avait ecrit la que la
%  relative postposee « جو ... ہے » devrait, une fois prise comme
%  outil ordinaire et non comme depannage, ramener les divergences a
%  deux par tableau au plus. Ce tableau, qui aligne quatre-vingt-dix
%  huit renvois et trente et une paires, en a rendu ZERO au premier
%  jet. Elle a servi quatre fois : la digue (4) avant le port (5), le
%  pont (48) avant le voilier (49), le plateau (96) avant le fort
%  (97), l'ile (99) avant les recifs (100). Il faudra deux tableaux
%  de plus avant d'en faire une regle, mais la mesure est posee.
%
%  TRENTE ET UNE PAIRES POUR TRENTE-SEPT BLOCS. Le tableau est un
%  INVENTAIRE presque pur — il nomme des navires, des grades, des
%  marchandises — et l'inventaire est ce qui coute le moins : les
%  renvois s'y suivent en enumeration, et l'enumeration ne se
%  retourne pas. Le tableau 9, qui decrivait, en demandait
%  trente-cinq pour dix blocs de moins.
% ===================================================================

%%K t10-apar-1 apar
\VUsaut{4.0mm}
\VUtitre{15.0pt}{60}{جدول نمبر 10}
\VUsaut{3.0mm}
\VUpk{11.4pt}{40}{سمندر۔ --- بندرگاہ۔}
\VUsaut{3.0mm}

%%K t10-01-1 p
1. --- گرمیوں میں کچھ لوگ پہاڑ پر سکون اور ٹھنڈک ڈھونڈتے ہیں، اور
کچھ ساحل کی طرف جانا پسند کرتے ہیں۔ ہم نے پچھلے سال یہی کیا، کیونکہ
ہمیں \VUgras{بحرِ اوقیانوس} \textsuperscript{(1)} بہت بھاتا ہے۔

%%K t10-02-1 p
2. --- جہاں تک میرا تعلق ہے، مجھے پانی کی اُس بے کراں وسعت کو دیکھنے
میں بڑا لطف آتا ہے جو \VUgras{افق} \textsuperscript{(2)} تک پھیلی ہوئی
ہے، اور اُن بہت بڑے \VUgras{بھاپ کے جہازوں} \textsuperscript{(3)} کو
لمبے \VUgras{سمندری سفر} پر روانہ ہوتے دیکھنے میں، جن پر امریکہ جانے
والے مسافر سوار ہوتے ہیں۔

%%K t10-03-1 p
3. --- اِسی لیے میرے والد، جو میرے شوق کو جانتے ہیں، چھٹیاں گزارنے
کے لیے ہمیشہ کوئی ایسا نہایت پُرسکون ساحل چنتے ہیں جو ہمارے کسی بڑے
\VUgras{بحری اڈّے} کے قریب ہو۔

%%K t10-03-2 p
ہمارے پچھلے قیام میں \VUgras{بحری بیڑا} اُس بہت بڑے \VUgras{بند}
\textsuperscript{(4)} کے پیچھے لنگر ڈالنے آیا جو \VUgras{فوجی
بندرگاہ} \textsuperscript{(5)} کو بند کرتا اور اُس کی حفاظت کرتا ہے،
اور میں جی بھر کر طرح طرح کے \VUgras{جنگی جہاز} دیکھ سکا — چھوٹی
\VUgras{آبدوز} \textsuperscript{(10)} بھی، جو ڈوب کر پانی کے نیچے غائب
ہو جاتی ہے، اور وہ بہت بڑا \VUgras{زرہ پوش جہاز}
\textsuperscript{(9)} بھی، جو اب تک تقریباً صرف \VUgras{تارپیڈو کشتی}
\textsuperscript{(8)} کے حملوں سے ڈرتا تھا۔

%%K t10-tit-2 sub
\VUsaut{3.0mm}
\VUpk{10.2pt}{40}{چھوٹے جہاز۔}
\VUsaut{3.0mm}

%%K t10-04-1 p
4. --- یہ کشتی پہلے ہی بڑی \VUgras{مہارت} سے موٹی دھاتی زرہ کا کمزور
مقام ڈھونڈ لینا جانتی تھی تاکہ وہاں خطرناک \VUgras{تارپیڈو} چپکا
سکے، جس کا دھماکہ جہاز کو ڈبو دیتا ہے ؛ پھر بھی وہ آبدوز سے کم ڈراؤنی
تھی، کیونکہ تارپیڈو شکن جہاز آسانی سے اُس کا پیچھا کر لیتے ہیں۔

%%K t10-05-1 p
5. --- میں تیز رفتار زرہ پوش \VUgras{کروزر} \textsuperscript{(6)} بھی
دیکھ سکا، جو اصل زرہ پوش جہاز کی طرح تقریباً ناقابلِ نقصان ہیں ؛ کئی
\VUgras{بار بردار جہاز} \textsuperscript{(12)}، تین چار
\VUgras{توپ کشتیاں} \textsuperscript{(7)} اور آخر میں ایک
\VUgras{فریگیٹ} \textsuperscript{(11)}۔ \VUgras{اُس بیڑے کا منظر، جب
وہ لنگر اٹھانے کے لیے چالیں چلتا ہے، سب سے دل چسپ ہوتا ہے۔}

%%K t10-06-1 p
6. --- اُن فولاد کے بڑے انباروں اور اُن خوش وضع \VUgras{سہ مستولی
جہازوں} یا اُن خوش نما \VUgras{بادبانی جہازوں} \textsuperscript{(13)}
کی ساخت میں کتنا فرق ہے، جو آج بھی تجارتی بیڑے میں کام آتے ہیں اور
جنھیں میں نے پورے بادبان کھولے بندرگاہ میں داخل ہوتے دیکھا، جب میں
\VUgras{جیٹی} \textsuperscript{(14)} پر ٹہل رہا تھا۔ سچ ہے کہ
\VUgras{ہوا} مخالف ہونے پر یہ اپنے بہت سے فائدے کھو بیٹھتے ہیں۔
اُنھیں حوض میں دوبارہ داخل ہونے کے لیے سارے بادبان اتارنے پڑتے ہیں،
اور اُنھیں لے جانے اور گھاٹ تک پہنچانے کے لیے ایک \VUgras{ٹگ}
\textsuperscript{(16)} درکار ہوتا ہے — جیسے وہ سادہ \VUgras{مال بردار
کشتیاں} \textsuperscript{(17)} ہوں۔

%%K t10-tit-3 sub
\VUsaut{3.0mm}
\VUpk{10.2pt}{40}{تجارتی بندرگاہ۔}
\VUsaut{3.0mm}

%%K t10-07-1 p
7. --- جس بندرگاہ کی میں بات کر رہا ہوں، اُس کی دل چسپی یہ ہے کہ اُس
میں صرف فوجی بندرگاہ ہی نہیں — جو دیو قامت کاموں سے بنائی گئی ہے —
بلکہ ایک حیرت انگیز \VUgras{تجارتی بندرگاہ} بھی ہے، جو ایک بڑے دریا
کے \VUgras{دہانے} پر واقع ہے۔

%%K t10-08-1 p
8. --- زمین کی جو پٹّی دونوں حوضوں کو جدا کرتی ہے، وہ مضبوط کی گئی
ہے اور \VUgras{توپ خانوں} اور \VUgras{برجوں} کے ایک سلسلے سے محفوظ
ہے، جو سمندر کی طرف بڑھتے ہیں۔ \VUgras{فصیلوں} \textsuperscript{(15)}
پر ٹہلنے کے لیے خاص اجازت درکار ہوتی ہے۔ وہ مجھے اپنے چچا کے ذریعے
آسانی سے مل گئی تھی، جو \VUgras{جہاز سازی کے کارخانوں}
\textsuperscript{(18)} کے انجینئر ہیں، اور میں وہ جہاز دیکھنے گیا جو
وسیع شیڈوں کے نیچے بنائے جاتے ہیں۔

%%K t10-09-1 p
9. --- میں تو ایک زرہ پوش جہاز کے پانی میں اتارے جانے پر بھی موجود
تھا، اور مجھے وہ دل ہلا دینے والا لمحہ یاد ہے جو سارے مجمعے نے محسوس
کیا، جب وہ بہت بڑا ڈھانچہ، اپنے حال پر چھوڑا ہوا، اپنی جھولے نما
چوکھٹ پر سرکنے لگا اور دریا کے پانی پر آ تیرا۔

%%K t10-10-1 p
10. --- کارخانوں تک جانے کے لیے \VUgras{مشقی میدان}
\textsuperscript{(19)} میں سے گزرنا پڑتا ہے، جہاں چھاؤنی کے سپاہی روز
مشق کرنے آتے ہیں ؛ پھر ایک لوہے کے \VUgras{چھوٹے پُل}
\textsuperscript{(20)} پر سے گزرتے ہیں، جو پریڈ کے میدان کو
\VUgras{اسلحہ خانے} \textsuperscript{(21)} سے ملاتا ہے۔

%%K t10-tit-4 sub
\VUsaut{3.0mm}
\VUpk{10.2pt}{40}{اسلحہ خانہ اور گودیاں۔}
\VUsaut{3.0mm}

%%K t10-11-1 p
11. --- اپنے چچا کے ساتھ میں نے وہ بڑی عمارت دیکھی، جو
\VUgras{توپوں} \textsuperscript{(22)}، \VUgras{گولوں}
\textsuperscript{(23)} اور \VUgras{ہاوٹزر کے گولوں} سے بھری ہے۔ میں نے
وہ \VUgras{دھات کا کارخانہ} \textsuperscript{(24)} بھی دیکھا، جس میں
بھاپ کے جہازوں کے \VUgras{بوائلر} \textsuperscript{(25)} بنائے جاتے
ہیں۔

%%K t10-12-1 p
12. --- بالکل پاس ہی \VUgras{گودیاں} \textsuperscript{(26)} ہیں —
مرمت کی گودیاں — اور وہ نہایت قابلِ ذکر \VUgras{تیرتی گودیاں}
\textsuperscript{(27)}، جن میں میں ایک زیرِ مرمت جہاز پر بہت قریب سے
اُس کا \VUgras{پروپیلر} \textsuperscript{(28)}، اُس کا \VUgras{تلا}
\textsuperscript{(29)} اور اُس کا \VUgras{پتوار} \textsuperscript{(30)}
دیکھ سکا۔

%%K t10-13-1 p
13. --- تجارتی بندرگاہ بھی فوجی بندرگاہ ہی کی طرح دیکھنے کے لائق
ہے، اور میں نے کئی گھنٹے اُن گھاٹوں پر منہ اٹھائے گزارے جہاں وہ
\VUgras{چرخی والے جہاز} \textsuperscript{(31)} بندھے رہتے ہیں جو دریا
کے اوپر کی طرف جاتے ہیں ؛ اور اُن \VUgras{مستول دار کرینوں}
\textsuperscript{(32)} کو کام کرتے دیکھتے ہوئے، جو بہت بھاری بوجھ پروں
کی طرح اٹھا لیتی ہیں۔

%%K t10-tit-5 sub
\VUsaut{4.0mm}
\VUpk{10.2pt}{40}{بندرگاہ کا رہبر۔}
\VUsaut{3.0mm}

%%K t10-14-1 p
14. --- میں نے اُن \VUgras{بحر پار جہازوں} \textsuperscript{(34)} کو
سراہا جو بندرگاہ سے نکل رہے تھے، اپنے \VUgras{جھنڈے}
\textsuperscript{(35)} کھلی ہوا میں لہراتے ہوئے، اور اُنھیں ایک ماہر
\VUgras{بندرگاہ کا رہبر} \textsuperscript{(36)} چلا رہا تھا، جو
\VUgras{تیرتے نشانوں} \textsuperscript{(40)} اور
\VUgras{نشان دار کھمبوں} سے نشان زد آبی گزرگاہ کے ساتھ ساتھ چلنا
حیرت انگیز طور پر جانتا ہے اور اُن \VUgras{چوڑی کشتیوں}
\textsuperscript{(41)} سے بچتا ہے جو اپنے اکلوتے چوکور بادبان سے
مشکل ہی سے سنبھلتی ہیں۔ میں نے \VUgras{اگلے حصے}
\textsuperscript{(37)} پر — یعنی جہاز کی نوک پر — دوسرے درجے کی
\VUgras{کیبنیں} \textsuperscript{(39)} پہچانیں، اور \VUgras{پچھلے
حصے} \textsuperscript{(38)} پر — یعنی دُم پر — پہلے درجے کے مسافروں
کے کمرے۔

%%K t10-15-1 p
15. --- کبھی کبھی میں ایسی \VUgras{مال بردار کشتی}
\textsuperscript{(42)} کی لدائی پر بھی موجود رہا جس میں ابھی ابھی
\VUgras{گودام} \textsuperscript{(43)} اور \VUgras{بحری اسٹیشن}
\textsuperscript{(44)} کا مال جمع کیا گیا تھا — وہ مال \VUgras{گھاٹ کی
ریل لائن} \textsuperscript{(45)} کی بے شمار گاڑیاں لاتی ہیں۔

%%K t10-tit-6 sub
\VUsaut{3.0mm}
\VUpk{10.2pt}{40}{چپّو کی سیر۔}
\VUsaut{3.0mm}

%%K t10-16-1 p
16. --- میں اکثر \VUgras{تیراؤ گودی} \textsuperscript{(53)} میں کشتی
چلاتا تھا، جس میں \VUgras{کشتیاں} \textsuperscript{(46)} پناہ لینے آتی
ہیں، اور \VUgras{چپّو} \textsuperscript{(47)} چلانا میرے لیے خوشی تھی۔
میں ایک \VUgras{عرشے} \textsuperscript{(48)} پر چڑھا جو ایک
\VUgras{بادبانی کشتی} \textsuperscript{(49)} کا تھا ؛
\VUgras{موٹے رسّوں} \textsuperscript{(52)} کے سہارے \VUgras{مستول}
\textsuperscript{(51)} پر چڑھ کر \VUgras{عرضی ڈنڈوں}
\textsuperscript{(50)} تک پہنچا ؛ پھر میں نے اپنی سیر
\VUgras{ڈونگی سے} \textsuperscript{(54)} دہرائی، بھاری
\VUgras{مچھیروں کی کشتیوں} \textsuperscript{(55)} کے بیچ، جہاں
\VUgras{جال} \textsuperscript{(56)} سوکھتے رہتے ہیں۔

%%K t10-17-1 p
17. --- \VUgras{گھاٹ} \textsuperscript{(58)} پر میرے سامنے سے طرح طرح
کا \VUgras{مال} \textsuperscript{(59)} گزرتا رہا، \VUgras{صندوقوں}
\textsuperscript{(60)} میں یا \VUgras{گٹھڑیوں} \textsuperscript{(61)}
میں بند۔ وہی بڑی کشتیوں کا \VUgras{بار} \textsuperscript{(63)} بنتا
تھا، اور زور دار \VUgras{کرینیں} \textsuperscript{(62)} اُسے اٹھا کر
\VUgras{لاریوں} \textsuperscript{(64)} پر رکھتی تھیں، جب کہ
\VUgras{مال بردار گاڑی بان} اُسے ظاہر کرتے اور \VUgras{محصول خانے}
\textsuperscript{(65)} میں محصول ادا کرتے تھے۔

%%K t10-18-1 p
18. --- آخر جب میں \VUgras{گھومنے والے پُل} \textsuperscript{(66)} یا
\VUgras{آبی پھاٹک} \textsuperscript{(69)} تک پہنچا، اُس \VUgras{گاد
نکالنے کی مشین} \textsuperscript{(68)} کے سامنے سے گزرتے ہوئے جو
\VUgras{نہر} \textsuperscript{(67)} کو صاف کرتی ہے، تو میں
\VUgras{اشارہ گاہ} \textsuperscript{(70)} کے قریب جیٹی پر اترا۔

%%K t10-19-1 p
19. --- اُس جیٹی کی دوسری طرف ہی وہ \VUgras{بڑی ناؤیں}
\textsuperscript{(71)} آ لگتی ہیں جو بیڑے کے افسروں کو خشکی پر لاتی
ہیں۔ ملّاحوں کو ایک تال پر اپنے چپّو اٹھاتے دیکھنا قابلِ تعریف منظر
ہے، جب کہ \VUgras{لہر} \textsuperscript{(72)} پر اٹھی ہوئی کشتی آسانی
سے اترنے کی سیڑھی سے جا لگتی ہے۔ اِسی جگہ سے \VUgras{مد و جزر} اور
\VUgras{ساحل} \textsuperscript{(73)} پر \VUgras{مد} اور \VUgras{جزر}
کی حرکت بہتر دیکھی جا سکتی ہے۔

%%K t10-tit-7 sub
\VUsaut{3.0mm}
\VUpk{10.2pt}{40}{کلب گھر۔}
\VUsaut{3.0mm}

%%K t10-20-1 p
20. --- گھر واپس آنے کے لیے میں نے \VUgras{بجلی کی ٹرام}
\textsuperscript{(74)} لی، جس کا \VUgras{اڈّا} \textsuperscript{(75)}
افسروں کے کلب کے سامنے گڑے \VUgras{کھمبے} کے پاس ہے۔

%%K t10-20-2 p
کلب کے اُس \VUgras{چھجّے} \textsuperscript{(76)} سے، جو \VUgras{لنگروں}
\textsuperscript{(77)}، توپوں اور گولوں سے سجا ہے، میں نے اکثر بحری
افسروں کا شان دار مجمع دیکھا : \VUgras{لیفٹیننٹ}
\textsuperscript{(78)} اپنی \VUgras{تلوار} کے ساتھ،
\VUgras{توپ خانے} \textsuperscript{(79)} اور \VUgras{پیادہ فوج}
\textsuperscript{(80)} کے کپتان اپنی \VUgras{شمشیر} کے ساتھ۔ اُن کے
پیچھے ایک \VUgras{ملّاح} \textsuperscript{(81)} تھا، جو اُنھیں
\VUgras{دوربین} لا دیتا تھا جب وہ دور کی چیز پہچاننا چاہتے تھے۔

%%K t10-21-1 p
21. --- میں نے جوان \VUgras{نائب لیفٹیننٹ} \textsuperscript{(82)} بھی
دیکھے، جو اپنے \VUgras{کمانڈر} \textsuperscript{(83)} یا
\VUgras{ایڈمرل} \textsuperscript{(84)} سے ہدایتیں لے رہے تھے، جب کہ
ایک \VUgras{کوارٹر ماسٹر} \textsuperscript{(85)} اُنھیں جہاز تک لے
جانے کے لیے انتظار کر رہا تھا۔

%%K t10-22-1 p
22. --- اُس چھجّے سے نظارہ شان دار ہے : سارا ساحل اپنے
\VUgras{خیموں} \textsuperscript{(86)} اور \VUgras{کیبنوں}
\textsuperscript{(87)} سمیت نیچے دکھائی دیتا ہے، اور نہانے کے وقت
\VUgras{نہانے والے} \textsuperscript{(88)} \VUgras{جوا خانے}
\textsuperscript{(89)} کے سامنے خوشی سے اچھلتے کودتے نظر آتے ہیں۔

%%K t10-23-1 p
23. --- ساحل کے سرے پر خشکی ایک چھوٹا \VUgras{چٹانی}
\VUgras{جزیرہ نما} \textsuperscript{(91)} بناتی ہے، جہاں \VUgras{موج}
\textsuperscript{(92)} آ کر ٹوٹتی ہے اور جس پر ایک \VUgras{روشنی کا
مینار} \textsuperscript{(93)} بنا ہے، جو رات کو جہاز رانوں کو راستہ
دکھاتا ہے۔ وہ جزیرہ نما خشکی سے صرف ایک بہت تنگ \VUgras{برزخ}
\textsuperscript{(90)} کے ذریعے ملا ہوا ہے۔

%%K t10-tit-8 sub
\VUsaut{3.0mm}
\VUpk{10.2pt}{40}{ساحلوں کا عمومی منظر۔}
\VUsaut{3.0mm}

%%K t10-24-1 p
24. --- \VUgras{کھڑی چٹان} \textsuperscript{(94)} دور تک پھیلی ہوئی
ہے، سمندر نے اُسے جگہ جگہ چیر رکھا ہے اور اپنی موج میں اُس میں
\VUgras{خلیجوں} \textsuperscript{(95)} اور \VUgras{راسوں} کا ایک
سلسلہ کھینچ دیا ہے، جس سے وہ نہایت دل کش ہو گئی ہے۔

%%K t10-24-2 p
\VUgras{سطح مرتفع} \textsuperscript{(96)} پر، جو \VUgras{آبنائے}
\textsuperscript{(98)} پر حاوی ہے، ایک بہت اہم \VUgras{قلعہ}
\textsuperscript{(97)} ہے اور چند « کوٹھیاں »۔

%%K t10-25-1 p
25. --- وہ آبنائے براعظم سے ایک نہایت خطرناک \VUgras{جزیرہ}
\textsuperscript{(99)} جدا کرتی ہے، جو \VUgras{زیرِ آب چٹانوں}
\textsuperscript{(100)} اور \VUgras{شکن موجوں} سے گھرا ہوا ہے، جہاں
کشتیاں اور بڑے جہاز تک کبھی کبھی جا لگتے ہیں۔ مدد کی تیاری اور
\VUgras{بچاؤ کشتیوں} کے جلد پہنچنے کے باوجود یہ \VUgras{جہاز کے
حادثے} دہشت ناک ہوتے ہیں، اور اعتدالِ شب و روز کے \VUgras{طوفانوں}
میں کئی جہاز آدمیوں اور مال سمیت تباہ ہو چکے ہیں۔

%%K t10-25-2 p
اِس ہمسائیگی کے باوجود \VUgras{بندرگاہ} پر ملّاحوں کا بہت آنا جانا
رہتا ہے۔

\VUsaut{6.0mm}
\VUfilet{22mm}
